\section{总结与心得体会}

本次课程设计以凯尔纳目镜为研究对象，采用Zemax OpticStudio软件完成了光学建模、仿真分析与优化设计的全过程。设计结果表明，通过合理设置评价函数并综合运用局部优化（DLS算法）与全局优化（锤形优化）等手段，最终得到了进一步优化后的凯尔纳目镜方案：有效焦距30.0 mm、最大径向视场22.5°、入瞳直径4 mm、后焦距21.78249 mm、系统总长46.19713 mm、最大畸变约1.7\%（小于5\%）、全视场相对照度大于85\%、全视场RMS波前误差约0.13\(\lambda\)（接近衍射极限），各项核心性能指标均满足设计要求。

通过公差分析进一步验证了设计的工程可行性，在合理加工公差范围内，约90\%以上的蒙特卡洛样本能够满足像质要求，通过设置补偿器可将良品率进一步提升至约92\%，具有良好的工程实用性。

关于目镜设计，尚有以下几点想法在本次设计中未能完全实现，可在未来的实践中尝试：

\begin{enumerate}
\def\labelenumi{\arabic{enumi}.}
\item
  玻璃材料系统性优化：本次仿真固定使用BK7和F2，未对玻璃材料进行筛选。如果以后有机会继续对该设计进行优化，可以多尝试几种材料组合，寻找在轻便性、透过率、成像质量和价格等方面均能满足要求的最优方案；
\item
  扩大视场角的尝试：通过引入四片甚至五片式目镜结构（如Plössl目镜、Erfle目镜），可以将表观视场角扩大至±25°甚至±35°，显著提升观测体验，但这需要更多的设计自由度和更复杂的优化策略；
\item
  非球面的引入：本次设计中所有镜面均为球面，如引入非球面可以更有效地校正高阶球差和彗差，进一步提升全视场的成像质量。Zemax支持多种非球面类型（偶次非球面、Q-型非球面等），可作为后续优化的探索方向；
\item
  高阶非球面系数的范围限制：经过本次仿真发现，非球面系数对偶次非球面的结构影响巨大，尤其是高阶系数。在今后的仿真中可以尝试对高阶非球面系数的范围加以限定，以达到更好的公差分析效果；
\item
  实物加工验证：若条件允许，可将优化后的设计方案提交实际加工，通过干涉仪测量面型误差和波前误差，验证仿真结果的准确性，并根据实测数据进行反馈修正。
\end{enumerate}

最后，感谢老师在应用光学及光学课程设计课堂上的悉心教导，使我对应用光学与光学设计有了更深入的认识，在上机时为我们耐心解答疑问和讨论，帮助我们顺利完成了本次设计。与同组同学的交流也让我较快地理解了初始结构选取与优化函数设置的关键要领，共同顺利完成了系统设计任务。
